\chapter{{\fontsize{14pt}{21pt}\selectfont\MakeUppercase{ПРОЕКТИРОВАНИЕ И ТЕХНИЧЕСКАЯ РЕАЛИЗАЦИЯ СИСТЕМЫ УПРАВЛЕНИЯ}}}
\label{cha:design}

% ============================================================
%  ГЛАВА 2. ПРОЕКТИРОВАНИЕ И ТЕХНИЧЕСКАЯ РЕАЛИЗАЦИЯ СИСТЕМЫ УПРАВЛЕНИЯ
% ============================================================
\chapter{Проектирование и техническая реализация системы управления}

\section{Разработка архитектуры распределённого программного комплекса}

\subsection{Выбор архитектурного стиля}

Проектирование архитектуры является фундаментальным этапом разработки программного комплекса: именно на данной стадии закладываются структурные свойства системы, определяющие её масштабируемость, надёжность, сопровождаемость и производительность в долгосрочной перспективе. В рамках настоящего раздела проводится анализ применимых архитектурных подходов, обосновывается выбор итогового архитектурного решения и формализуется структура разработанного программного комплекса.

В современной инженерии программного обеспечения наибольшее распространение получили два конкурирующих архитектурных стиля: монолитная и микросервисная архитектуры. Для обоснованного выбора необходимо рассмотреть их характеристики в контексте требований, сформулированных в первой главе настоящей работы.

Микросервисная архитектура предполагает декомпозицию системы на множество независимо разворачиваемых сервисов, каждый из которых реализует конкретную бизнес-функцию и взаимодействует с остальными посредством сетевых протоколов --- как правило, HTTP REST или gRPC. Данный подход демонстрирует преимущества в условиях работы крупных распределённых команд разработчиков, при необходимости независимого масштабирования отдельных компонентов и при высокой частоте выпуска обновлений изолированных частей системы. К характерным достоинствам микросервисной архитектуры относятся: независимость технологического стека отдельных сервисов, устойчивость к отказам на уровне изолированных процессов, а также возможность горизонтального масштабирования конкретных узлов без перезапуска всей системы.

Вместе с тем микросервисный подход сопряжён с существенными эксплуатационными издержками, принципиально важными в контексте настоящего проекта. Во-первых, каждый сервис требует собственной инфраструктуры развёртывания: отдельный сетевой порт, конфигурационный файл, процесс инициализации и мониторинга. Во-вторых, взаимодействие сервисов по сети порождает дополнительные задержки и накладные расходы на сериализацию и десериализацию данных при каждом межсервисном вызове. В-третьих, полноценная реализация микросервисной архитектуры требует развёртывания дополнительных инфраструктурных компонентов: системы оркестрации контейнеров, сервисного реестра и системы трассировки запросов. Перечисленные компоненты существенно увеличивают сложность системы и создают зависимости от зарубежного программного обеспечения, что противоречит требованию технологической независимости, сформулированному в первой главе.

Монолитная архитектура объединяет все компоненты системы в рамках единого процесса операционной системы, разделяя общее адресное пространство памяти, единую точку входа и общую конфигурацию развёртывания. Несмотря на устойчивое представление о монолите как об устаревшем архитектурном решении, данный подход демонстрирует высокую эффективность в условиях, характерных для разрабатываемого комплекса: команда разработки состоит из одного специалиста, функциональные требования чётко определены и стабильны, а потребности в раздельном масштабировании отдельных внутренних компонентов отсутствуют.

Ключевым аргументом в пользу монолитного подхода является специфика решаемой задачи: разрабатываемая система является «тонким» координирующим слоем над управляемыми имитационными сервисами. Её основная функция --- оркестрация процессов операционной системы на целевых хостах, маршрутизация входящего трафика и хранение конфигурационных данных. Введение внутрисистемных сетевых вызовов в рамках гипотетической микросервисной декомпозиции создало бы накладные расходы, несоразмерные выигрышу от разделения компонентов.

Распределённость системы, заявленная в названии раздела, реализуется не на уровне внутреннего устройства приложения, а на уровне топологии развёртывания: единый управляющий процесс координирует целевые хосты посредством защищённых сетевых протоколов. Таким образом, распределённость является внешней характеристикой инфраструктуры, а не принципом внутренней декомпозиции.

На основании проведённого анализа для реализации программного комплекса выбрана монолитная архитектура с явно выраженной внутренней модульной структурой. Данное решение обеспечивает минимальные накладные расходы на взаимодействие компонентов, простоту развёртывания и эксплуатации в закрытом корпоративном контуре, технологическую независимость от внешних оркестраторов, а также высокую сопровождаемость и читаемость единой кодовой базы.

В таблице~\ref{tab:arch-compare} приведён сравнительный анализ двух рассмотренных архитектурных подходов по критериям, наиболее значимым для данного проекта.

\begin{table}[H]
\caption{Сравнение архитектурных подходов}
\label{tab:arch-compare}
\centering
\begin{tabular}{|p{5cm}|p{4.5cm}|p{4.5cm}|}
\hline
Критерий & Монолитная архитектура & Микросервисная архитектура \\
\hline
Накладные расходы на взаимодействие компонентов & Минимальные (вызовы функций) & Высокие (сетевые вызовы, сериализация) \\
\hline
Сложность развёртывания & Низкая (один процесс) & Высокая (множество сервисов, оркестратор) \\
\hline
Зависимость от внешних систем & Отсутствует & Требует оркестратора (Kubernetes и др.) \\
\hline
Сопровождаемость при одном разработчике & Высокая & Низкая \\
\hline
Независимое масштабирование компонентов & Не поддерживается & Поддерживается \\
\hline
\end{tabular}
\end{table}

\subsection{Топология развёртывания и гибкость инфраструктуры}

Одним из ключевых проектных решений, определивших структуру системы, является принцип гибкого размещения имитационных сервисов. В отличие от систем, жёстко делящихся на «локальный» и «кластерный» режимы, разработанный программный комплекс предоставляет оператору возможность назначать произвольный целевой хост для каждой заглушки независимо.

Распределённость системы реализуется посредством реестра хостов: каждый зарегистрированный хост описывается своим сетевым адресом, учётной записью для подключения и рабочей директорией. При развёртывании конкретной заглушки оператор выбирает один из хостов реестра в качестве целевого. Это позволяет в рамках единого экземпляра управляющего приложения одновременно:

\begin{itemize}
\item размещать часть заглушек непосредственно на сервере, где работает управляющее приложение (адрес \texttt{127.0.0.1} или hostname сервера), --- что соответствует традиционному «локальному» сценарию;
\item размещать другие заглушки на удалённых серверах через SSH/SCP --- что соответствует «кластерному» сценарию;
\item совмещать оба варианта в рамках одной инфраструктуры без каких-либо изменений в конфигурации приложения.
\end{itemize}

Таким образом, граница между «локальным» и «распределённым» режимами определяется не конфигурацией запуска приложения, а составом реестра хостов и выбором оператора при регистрации каждой заглушки. Такой подход значительно увеличивает гибкость эксплуатации системы и позволяет плавно наращивать инфраструктуру без вмешательства в работу уже запущенных сервисов.

С точки зрения внутренней архитектуры это означает, что модуль управления жизненным циклом всегда работает по единому алгоритму, независимо от того, является ли целевой хост локальным или удалённым. При работе с \texttt{localhost} SSH-соединение устанавливается к локальному SSH-серверу по стандартному порту 22, что унифицирует код и устраняет необходимость в двух ветках логики для локального и удалённого сценариев.

\subsection{Контекстная диаграмма системы}

Для формализованного описания архитектуры программного комплекса применяется нотация C4~Model, предусматривающая четыре уровня детализации: контекстный (Level~1), контейнерный (Level~2), компонентный (Level~3) и кодовый (Level~4). Нотация обеспечивает однозначность интерпретации архитектурных решений и соответствует требованиям методических указаний по применению стандартизованных схем для представления структуры системы.

Контекстная диаграмма (рисунок~\ref{fig:c4-context}) отображает систему как единый «чёрный ящик» в окружении внешних акторов и смежных систем. На данном уровне абстракции детали внутреннего устройства системы не раскрываются; основное внимание уделяется определению границ системы и характеру её взаимодействия с внешним окружением.

\begin{figure}[H]
\centering
\includegraphics[width=\textwidth]{SystemContext.png}
\caption{Контекстная диаграмма системы (C4 Level~1)}
\label{fig:c4-context}
\end{figure}


% workspace {

%     model {

%         engineer = person "Инженер тестирования" {
%             description "Управляет инфраструктурой заглушек через веб-интерфейс"
%         }

%         consumer = person "Потребитель" {
%             description "Любой клиент, направляющий HTTP-запросы к заглушкам: нагрузочный инструмент, тестируемая система или иной HTTP-клиент"
%         }

%         prometheus = softwareSystem "Prometheus" {
%             description "Внешняя система сбора и хранения метрик"
%         }

%         mockSystem = softwareSystem "Программный комплекс управления заглушками" {
%             description "Монолитное Python/FastAPI-приложение"

%             webUI = container "Веб-интерфейс (Jinja2 SSR)" {
%                 description "HTML-страницы, формируемые на стороне сервера: Dashboard, Настройки, Логи"
%                 technology "FastAPI / Jinja2 / HTML"
%             }

%             apiLayer = container "REST API" {
%                 description "Бизнес-логика: управление заглушками, проксирование, настройки"
%                 technology "Python 3.10+ / FastAPI / Uvicorn (ASGI)"

%                 lifecycleMgr = component "Lifecycle Manager" {
%                     description "Регистрация (загрузка + SCP на хост), запуск, остановка, удаление Java-процессов"
%                     technology "Python asyncio / paramiko (SSH/SCP)"
%                 }

%                 proxyEngine = component "Proxy Engine" {
%                     description "Прозрачная маршрутизация HTTP-запросов от потребителей к целевым заглушкам"
%                     technology "Python httpx (async)"
%                 }

%                 rateLimiter = component "Rate Limiter" {
%                     description "Алгоритм Fixed Window Counter: ограничение RPS на уровне каждой заглушки"
%                     technology "Redis INCR + EXPIRE"
%                 }

%                 settingsModule = component "Settings Module" {
%                     description "Управление учётными записями SSH и реестром хостов развёртывания"
%                     technology "Python / Redis"
%                 }

%                 cryptoService = component "Crypto Service" {
%                     description "Шифрование/дешифрование паролей SSH перед сохранением в Redis"
%                     technology "Python cryptography (Fernet)"
%                 }

%                 metricsExporter = component "Metrics Exporter" {
%                     description "Эндпоинт /metrics в формате Prometheus exposition format"
%                     technology "Python"
%                 }
%             }

%             redis = container "Redis" {
%                 description "In-Memory хранилище конфигураций и состояний"
%                 technology "Redis 7.x (AOF persistence)"
%             }
%         }

%         targetHost = softwareSystem "Целевой хост" {
%             description "Сервер ОС Astra Linux с запущенными заглушками. Может совпадать с сервером управляющего приложения."
%         }

%         engineer   -> webUI        "Управляет через браузер (HTTP)"
%         consumer   -> proxyEngine  "HTTP-запросы к заглушкам"
%         prometheus -> apiLayer     "Scrape /metrics (HTTP GET)"

%         webUI          -> apiLayer     "Вызовы REST API (HTTP)"
%         apiLayer       -> redis        "Чтение/запись конфигураций и состояний"
%         lifecycleMgr   -> targetHost   "SCP артефакта при регистрации; SSH: запуск/остановка/удаление процесса и файла"
%         proxyEngine    -> targetHost   "Перенаправление HTTP-запросов к работающей заглушке"
%         settingsModule -> cryptoService "Шифрование пароля перед сохранением; дешифрование перед использованием"
%         rateLimiter    -> redis        "Инкремент счётчика и проверка лимита"
%     }

%     views {

%         systemContext mockSystem "SystemContext" {
%             include *
%             autolayout lr
%             title "Рисунок 2.1 – Контекстная диаграмма системы (C4 Level 1)"
%         }

%         container mockSystem "Containers" {
%             include *
%             autolayout lr
%             title "Рисунок 2.2 – Контейнерная диаграмма (C4 Level 2)"
%         }

%         component apiLayer "Components" {
%             include *
%             autolayout tb
%             title "Рисунок 2.3 – Компонентная диаграмма REST API (C4 Level 3)"
%         }

%         styles {
%             element "Person" { 
%             shape Person  
%             background #08427B 
%             color #ffffff 
                
%             }
%             element "Software System" { 
%             background #1168BD 
%             color #ffffff 
                
%             }
%             element "Container"       { 
%             background #438DD5 
%             color #ffffff 
                
%             }
%             element "Component"       { 
%             background #85BBF0 
%             color #000000 
                
%             }
%         }
%     }
% }



Внешними акторами системы являются:

\begin{itemize}
\item инженер тестирования --- специалист, осуществляющий управление инфраструктурой имитационных сервисов через веб-интерфейс системы; выполняет операции регистрации артефактов, запуска и остановки заглушек, настройки параметров и просмотра метрик;
\item потребитель --- любой HTTP-клиент, направляющий запросы к имитационным сервисам через прокси-движок системы: это может быть инструмент нагрузочного тестирования (Apache JMeter, Gatling и др.), тестируемая система (System Under Test) или иной HTTP-клиент;
\item Prometheus --- внешняя система сбора и хранения метрик, периодически опрашивающая специализированный эндпоинт системы для получения телеметрических данных;
\item целевой хост --- сервер под управлением ОС Astra Linux, на котором непосредственно выполняются Java-процессы имитационных сервисов; может совпадать с сервером управляющего приложения.
\end{itemize}

\subsection{Контейнерная диаграмма}

Контейнерная диаграмма (рисунок~\ref{fig:c4-container}) раскрывает внутреннее устройство программного комплекса на уровне самостоятельно исполняемых программных единиц --- контейнеров. В терминологии C4~Model «контейнер» обозначает логически обособленную программную единицу, а не технологию Docker-контейнеризации. Для каждого контейнера указывается его технологический стек и зона ответственности.

\begin{figure}[H]
\centering
\includegraphics[width=\textwidth]{Containers.png}
\caption{Контейнерная диаграмма системы (C4 Level~2)}
\label{fig:c4-container}
\end{figure}

В составе программного комплекса выделяются следующие контейнеры.

Веб-интерфейс (Jinja2 SSR). Реализует пользовательский интерфейс на основе технологии Server-Side Rendering с использованием шаблонизатора Jinja2, встроенного в экосистему FastAPI. HTML-страницы формируются на стороне сервера и передаются браузеру в готовом виде. Данный подход минимизирует объём JavaScript-кода на стороне клиента и обеспечивает полную функциональность интерфейса без зависимости от внешних CDN. Веб-интерфейс включает: главную панель (Dashboard) со списком заглушек, страницу управления хостами и настроек системы, страницу просмотра журналов процессов.

REST API (FastAPI). Центральный контейнер системы, реализующий всю бизнес-логику. Запускается под управлением ASGI-сервера Uvicorn. REST API предоставляет конечные точки для взаимодействия с веб-интерфейсом, а также программный интерфейс для возможной интеграции со сторонними системами автоматизации. Внутренняя структура контейнера детализируется на компонентном уровне в разделе~2.1.5.

Redis. In-Memory хранилище структур данных, выполняющее функции постоянного хранилища конфигурационных данных системы и оперативного хранилища счётчиков алгоритма ограничения запросов. Redis функционирует в режиме персистентности AOF (Append Only File), что обеспечивает сохранность данных при аварийном перезапуске управляющей системы. Детальное описание структур данных Redis приводится в разделе~2.2.

Система намеренно не содержит выделенного контейнера хранилища артефактов. Архитектурное решение «1~заглушка --- 1~хост» исключает необходимость повторного развёртывания одного артефакта на несколько хостов, а значит, постоянное хранение загруженного файла на сервере управляющего приложения избыточно. При регистрации заглушки её бинарный файл принимается из браузера во временный каталог, немедленно копируется на целевой хост посредством SCP и удаляется с сервера управляющего приложения. Единственной постоянной копией артефакта остаётся файл на целевом хосте в рабочей директории, указанной в реестре хостов.

\subsection{Компонентная диаграмма REST API}

Компонентная диаграмма (рисунок~\ref{fig:c4-component}) детализирует внутреннее устройство контейнера REST API, раскрывая составляющие его программные модули --- компоненты.

\begin{figure}[H]
\centering
\includegraphics[width=\textwidth]{Components.png}
\caption{Компонентная диаграмма REST API (C4 Level~3)}
\label{fig:c4-component}
\end{figure}

\textbf{Lifecycle Manager.} Реализует операции управления жизненным циклом имитационных сервисов. При регистрации новой заглушки модуль принимает загруженный бинарный файл во временный каталог на сервере управляющего приложения, копирует его на целевой хост посредством SCP в рабочую директорию, указанную в записи хоста, и немедленно удаляет временный файл. Путь к файлу на целевом хосте не сохраняется в Redis --- он вычисляется при необходимости как конкатенация поля \texttt{working\_dir} из записи \texttt{hosts:\{hostname\}} и имени файла заглушки. При запуске модуль устанавливает SSH-соединение с целевым хостом, определяет свободный порт в заданном диапазоне и инициирует Java-процесс с указанными JVM-аргументами в фоновом режиме (\texttt{nohup}). При остановке выполняется корректное завершение процесса через SSH. При удалении заглушки производятся: остановка процесса (если запущен), удаление бинарного файла с целевого хоста командой \texttt{SSH rm} и удаление записи из Redis.

\textbf{Proxy Engine.} Реализует функциональность прозрачного обратного прокси-сервера на прикладном уровне. При поступлении HTTP-запроса движок извлекает адрес и порт целевой заглушки из Redis, формирует новый HTTP-запрос с сохранением метода, заголовков и тела оригинального запроса, а затем возвращает ответ заглушки инициатору в неизменном виде. Для асинхронного выполнения исходящих запросов используется библиотека \texttt{httpx}.

\textbf{Rate Limiter.} Реализует алгоритм «Счётчик с фиксированным временным окном» (Fixed Window Counter). Компонент функционирует на уровне каждой заглушки независимо и управляется флагом \texttt{rate\_limit\_enabled}, переключаемым оператором через Dashboard без остановки заглушки. Детальное описание алгоритма приводится в разделе~2.4.

\textbf{Settings Module.} Обеспечивает управление конфигурационными параметрами системы --- учётными записями ОС Linux и реестром хостов развёртывания. Подробное описание приведено в разделе~2.1.6.

\textbf{Crypto Service.} Выполняет шифрование паролей SSH перед записью в Redis и дешифрование перед использованием для установки соединения. Применяется алгоритм AES-128 (Fernet). Ключ шифрования хранится на файловой системе сервера в файле с правами 600.

\textbf{Metrics Exporter.} Реализует эндпоинт \texttt{/metrics} в формате Prometheus exposition format. Агрегирует данные о RPS, срабатываниях Rate Limiter, статусах процессов и доступности хостов. Детальное описание приведено в разделе~2.6.

\subsection{Архитектура модуля настроек системы}

Модуль настроек (Settings Module) предоставляет администратору системы веб-интерфейс для конфигурирования двух категорий параметров: учётных записей операционной системы Linux, применяемых для SSH/SCP-операций, и реестра хостов, доступных для развёртывания имитационных сервисов.

\textbf{Управление учётными записями ОС Linux.} Данная категория позволяет зарегистрировать именованные учётные записи пользователей операционной системы Linux. Каждая учётная запись используется для аутентификации при установке SSH-соединения с целевым хостом и включает имя пользователя (login) и пароль. Применение отдельных учётных записей с ограниченными правами доступа обусловлено требованиями информационной безопасности: рабочая директория заглушек должна быть доступна только указанному пользователю, что минимизирует риски при нештатном поведении Java-процессов.

\textbf{Механизм защищённого хранения паролей.} Поскольку пароли учётных записей SSH хранятся в Redis, необходимо обеспечить их конфиденциальность в случае несанкционированного доступа к базе данных. Для этого применяется симметричное шифрование на основе алгоритма AES-128 в режиме CBC, реализованного в библиотеке \texttt{cryptography} пакетом Fernet. Схема работы следующая:

\begin{itemize}
\item при первом запуске приложения генерируется или загружается из защищённого конфигурационного файла ключ шифрования (Fernet key), который хранится вне Redis --- в файле с правами доступа 600, доступном только пользователю, от имени которого запущено приложение;
\item перед сохранением пароля в Redis компонент Crypto Service шифрует его с использованием Fernet-ключа и сохраняет в виде строки Base64-кодированного зашифрованного блока;
\item перед использованием пароля для установки SSH-соединения Crypto Service выполняет дешифрование, получая исходный пароль в открытом виде только в оперативной памяти процесса на время установки соединения.
\end{itemize}

Такой подход гарантирует, что даже при получении злоумышленником дампа Redis пароли останутся недоступными без Fernet-ключа, хранящегося на файловой системе сервера.

\textbf{Реестр хостов развёртывания.} Данная категория настроек представляет собой список серверов, доступных для развёртывания имитационных сервисов. Каждая запись реестра описывает конкретный целевой хост. В таблице~\ref{tab:host-params} приведено описание параметров модуля настроек с указанием типов данных и ограничений.

\begin{table}[H]
\caption{Параметры модуля настроек системы}
\label{tab:host-params}
\centering
\begin{tabular}{|p{3.5cm}|p{2.5cm}|p{7cm}|}
\hline
Параметр & Тип & Описание \\
\hline
\texttt{hostname} & String & Сетевое имя или IP-адрес целевого хоста; является уникальным идентификатором записи \\
\hline
\texttt{ssh\_port} & Integer & Порт SSH-сервера на целевом хосте (по умолчанию 22) \\
\hline
\texttt{account\_uuid} & UUID & Идентификатор учётной записи SSH из реестра \texttt{accounts} \\
\hline
\texttt{working\_dir} & String & Абсолютный путь к рабочей директории на целевом хосте для хранения файлов заглушек \\
\hline
\texttt{java\_path} & String & Абсолютный путь к исполняемому файлу Java на целевом хосте \\
\hline
\texttt{mock\_port\_min} & Integer & Нижняя граница диапазона TCP-портов для динамического назначения заглушкам \\
\hline
\texttt{mock\_port\_max} & Integer & Верхняя граница диапазона TCP-портов для динамического назначения заглушкам \\
\hline
\texttt{description} & String & Произвольное текстовое описание хоста для документирования в интерфейсе \\
\hline
\end{tabular}
\end{table}

Параметр \texttt{java\_path} введён в реестр хостов по причине неоднородности целевой инфраструктуры: в корпоративных средах переменная окружения \texttt{JAVA\_HOME} нередко не определена или указывает на версию JRE, не совместимую с требованиями конкретной заглушки. Указание явного пути к исполняемому файлу Java устраняет неоднозначность и гарантирует запуск процесса с нужной версией JDK вне зависимости от конфигурации окружения на хосте.

Параметры \texttt{mock\_port\_min} и \texttt{mock\_port\_max} задают диапазон TCP-портов, в пределах которого система будет динамически выбирать свободный порт при запуске каждой новой заглушки на данном хосте. Ограничение диапазона необходимо для совместимости с политиками межсетевого экрана, которые в закрытых корпоративных сетях, как правило, разрешают входящий трафик только на заранее согласованные порты.

Поле \texttt{description} позволяет операторам документировать назначение каждого хоста непосредственно в интерфейсе системы, что упрощает эксплуатацию при работе с несколькими окружениями.

\subsection{Диаграмма развёртывания}

Диаграмма развёртывания (рисунок~\ref{fig:deployment}) описывает физическое размещение компонентов системы на серверной инфраструктуре. Сервер управляющего приложения содержит два постоянных компонента: FastAPI-процесс под управлением Uvicorn и экземпляр Redis. Бинарные файлы заглушек постоянно хранятся исключительно на целевых хостах в рабочих директориях. Временный каталог на сервере управляющего приложения используется лишь на время передачи файла (в течение одной SCP-операции) и очищается немедленно после её завершения.

\begin{figure}[H]
\centering
% \includegraphics[width=\textwidth]{images/deployment.png}
\caption{Диаграмма развёртывания}
\label{fig:deployment}
\end{figure}

% @startuml deployment_diagram
% skinparam backgroundColor #FFFFFF
% skinparam defaultFontName Arial
% skinparam defaultFontSize 12

% skinparam node {
%     BackgroundColor #D6E4F0
%     BorderColor #1A6FA8
%     BorderThickness 2
% }

% skinparam rectangle {
%     BackgroundColor #FAFAFA
%     BorderColor #5D6D7E
% }

% node "**Сервер управляющего приложения**\n//Astra Linux SE 1.7 · Python 3.10 · Uvicorn · Redis 7.x//" as Master {
%     rectangle "**«FastAPI-процесс (Uvicorn)»**\nREST API + Веб-интерфейс\n+ все модули системы" as FastAPI
%     rectangle "**«Redis»**\nIn-Memory СУБД\nAOF-персистентность" as Redis
%     FastAPI --> Redis : Чтение/запись\nсостояний (TCP 6379)
% }

% node "**«localhost»**\n//Тот же физический сервер · OpenJDK 17//" as LocalHost #D5F5E3 {
%     rectangle "**«Java-процессы заглушек»**\nОдин или несколько\n.jar/.war процессов" as LocalMocks #FAFAFA
% }

% node "**«Удалённый хост 1»**\n//Astra Linux SE 1.7 · SSH · OpenJDK 17//" as Remote1 #FEF9E7 {
%     rectangle "**«Java-процессы заглушек»**\nОдин или несколько\n.jar/.war процессов" as RemoteMocks1 #FAFAFA
% }

% node "**«Удалённый хост 2»**\n//Astra Linux SE 1.7 · SSH · OpenJDK 17//" as Remote2 #FEF9E7 {
%     rectangle "**«Java-процессы заглушек»**\nОдин или несколько\n.jar/.war процессов" as RemoteMocks2 #FAFAFA
% }

% FastAPI --> LocalMocks   : Управление процессами (localhost)
% FastAPI --> RemoteMocks1 : SSH/SCP при регистрации;\nHTTP-проксирование запросов
% FastAPI --> RemoteMocks2 : SSH/SCP при регистрации;\nHTTP-проксирование запросов

% @enduml



На целевых хостах не требуется установки каких-либо компонентов разрабатываемой системы. Обязательные условия для включения хоста в реестр: работающий SSH-сервер, настроенная учётная запись пользователя с правами на запись в рабочую директорию и наличие исполняемого файла Java по пути, указанному в параметре \texttt{java\_path}. Бинарные файлы заглушек хранятся в рабочей директории хоста на протяжении всего времени существования записи о заглушке в системе и удаляются только при явном удалении заглушки оператором.

\subsection{Сценарии взаимодействия компонентов}

Для иллюстрации сквозного взаимодействия компонентов рассмотрим пять типовых сценариев. Принципиальное архитектурное решение: копирование бинарного артефакта на целевой хост выполняется однократно при регистрации заглушки; все последующие операции (запуск, остановка, переключение Rate Limiter) работают с файлом, уже находящимся на целевом хосте.

\textbf{Сценарий~1. Регистрация новой заглушки.}

\begin{enumerate}
\item Инженер тестирования загружает \texttt{.jar}- или \texttt{.war}-файл через веб-интерфейс и задаёт параметры: JVM-аргументы, целевой хост из реестра, пороговое значение Rate Limiter. Флаг \texttt{rate\_limit\_enabled} устанавливается в \texttt{false} по умолчанию.
\item REST API сохраняет файл во временный каталог на сервере управляющего приложения.
\item Lifecycle Manager получает параметры целевого хоста из Redis (адрес, учётная запись, рабочая директория). Crypto Service дешифрует пароль учётной записи.
\item Lifecycle Manager копирует файл с сервера управляющего приложения на целевой хост посредством SCP в рабочую директорию.
\item После успешного завершения SCP временный файл удаляется с сервера управляющего приложения.
\item В Redis создаётся JSON-запись заглушки с ключом \texttt{mocks:\{filename\}}. Имя файла добавляется в реестр \texttt{mocks:registry}. Статус устанавливается в \texttt{REGISTERED}.
\item Веб-интерфейс отображает новую заглушку на Dashboard со статусом «Зарегистрирована» и переключателем Rate Limiter в положении «Выкл.».
\end{enumerate}

\textbf{Сценарий~2. Запуск зарегистрированной заглушки.}

\begin{enumerate}
\item Инженер тестирования нажимает кнопку «Запустить» для выбранной заглушки на Dashboard.
\item Браузер отправляет POST-запрос к эндпоинту \texttt{/api/v1/mocks/\{mock\_name\}/start}.
\item Lifecycle Manager считывает из Redis конфигурацию заглушки (JVM-аргументы) и параметры хоста (адрес, \texttt{java\_path}, учётная запись). Crypto Service дешифрует пароль.
\item Lifecycle Manager определяет свободный TCP-порт в диапазоне \texttt{mock\_port\_min} / \texttt{mock\_port\_max} через SSH-команду на целевом хосте.
\item Lifecycle Manager выполняет SSH-команду: \texttt{\{java\_path\} \{jvm\_args\} -jar \{artifact\_path\} --server.port=\{port\}}. Процесс запускается в фоновом режиме (\texttt{nohup}).
\item PID процесса, назначенный порт и статус \texttt{RUNNING} сохраняются в запись заглушки в Redis.
\item Proxy Engine начинает маршрутизировать входящие запросы на зафиксированный порт заглушки.
\end{enumerate}

\textbf{Сценарий~3. Обработка входящего запроса через прокси.}

\begin{enumerate}
\item Потребитель направляет HTTP-запрос на адрес вида \texttt{/\{mock\_name\}/\{path\}}.
\item Proxy Engine извлекает из Redis запись заглушки по \texttt{mock\_name}: адрес, порт, статус, флаг \texttt{rate\_limit\_enabled}.
\item Если статус отличен от \texttt{RUNNING} --- возвращается HTTP~503 (Service Unavailable).
\item Если \texttt{rate\_limit\_enabled = false} --- запрос немедленно передаётся целевой заглушке, счётчики не обновляются. Если \texttt{rate\_limit\_enabled = true} --- Rate Limiter инкрементирует счётчик \texttt{rate:\{mock\_name\}:\{window\}} и сравнивает с порогом \texttt{rate\_limit}. При превышении --- возвращается HTTP~429 (Too Many Requests).
\item Proxy Engine асинхронно передаёт запрос целевой заглушке через \texttt{httpx} и возвращает ответ потребителю в неизменном виде.
\end{enumerate}

\textbf{Сценарий~4. Переключение Rate Limiter на Dashboard.}

\begin{enumerate}
\item Оператор переключает тумблер Rate Limiter для выбранной заглушки на Dashboard.
\item Браузер отправляет PATCH-запрос к \texttt{/api/v1/mocks/\{mock\_name\}/rate-limit} с телом \texttt{\{"enabled": true/false\}}.
\item Settings Module обновляет поле \texttt{rate\_limit\_enabled} в JSON-записи заглушки в Redis последовательностью операций GET~→~modify~→~SET. Для предотвращения состояния гонки при конкурентных изменениях операция выполняется с использованием механизма оптимистичной блокировки Redis (WATCH/MULTI/EXEC).
\item Начиная со следующего запроса к данной заглушке, Proxy Engine применяет обновлённое значение флага без перезапуска.
\end{enumerate}

\textbf{Сценарий~5. Удаление заглушки.}

\begin{enumerate}
\item Инженер тестирования инициирует удаление заглушки через веб-интерфейс.
\item Если заглушка в состоянии \texttt{RUNNING} --- Lifecycle Manager выполняет остановку процесса через SSH (SIGTERM, при необходимости SIGKILL).
\item Lifecycle Manager удаляет бинарный файл заглушки с целевого хоста командой \texttt{SSH rm \{artifact\_path\}}.
\item Запись о заглушке и все связанные ключи (счётчики Rate Limiter) удаляются из Redis командами DEL и SREM.
\item Запись исчезает из Dashboard.
\end{enumerate}

\section{Проектирование структур данных для хранения конфигураций и состояний в Redis}

\subsection{Обоснование модели данных}

Проектирование структур данных является фундаментальным этапом, определяющим производительность, корректность и сопровождаемость системы в целом. В разрабатываемом программном комплексе Redis выполняет роль единственного хранилища всех оперативных данных: конфигураций заглушек, параметров хостов и учётных записей, состояний запущенных процессов, счётчиков алгоритма Rate Limiter и вспомогательных индексов. Бинарные файлы артефактов в Redis не хранятся и не кэшируются --- единственной постоянной копией каждого артефакта является файл на целевом хосте.

Выбор Redis в качестве хранилища обусловлен его архитектурными характеристиками, подробно рассмотренными в разделе~1.3: время доступа к любому ключу составляет менее 1~мс, что критически важно для алгоритма Rate Limiter, выполняемого при обработке каждого входящего запроса.

Redis предоставляет несколько базовых типов данных: строки (String), хеши (Hash), списки (List), множества (Set) и сортированные множества (Sorted Set). Для хранения конфигурационных объектов --- заглушек, хостов, учётных записей --- выбран тип String в сочетании с JSON-сериализацией. Такой подход позволяет атомарно читать и записывать всю конфигурацию объекта одной операцией GET/SET, избегая частичных обновлений и связанных с ними состояний гонки (race conditions). Для хранения реестров (множеств идентификаторов) используется тип Set, гарантирующий уникальность элементов. Для счётчиков Rate Limiter применяется тип String со встроенными атомарными командами INCR и EXPIRE.

В таблице~\ref{tab:redis-namespaces} приведена сводная схема всех пространств имён ключей Redis, используемых в системе.



\subsection{Структура данных заглушки}

Центральным объектом системы является запись о заглушке --- имитационном Java-сервисе, управляемом программным комплексом. Запись хранится по ключу вида \texttt{mocks:\{filename\}}, где \texttt{\{filename\}} --- имя загруженного файла артефакта (например, \texttt{payment-service-stub.jar}). Имя файла является уникальным идентификатором заглушки в системе: при попытке зарегистрировать файл с именем, уже присутствующим в реестре \texttt{mocks:registry}, система возвращает ошибку и предлагает оператору переименовать загружаемый артефакт или удалить существующую запись.

Запись заглушки не хранит путь к бинарному файлу явно. Поскольку файл всегда располагается в рабочей директории хоста, путь вычисляется детерминировано при необходимости как \texttt{\{working\_dir\}/\{filename\}}, где \texttt{working\_dir} берётся из записи хоста, а \texttt{filename} --- из ключа Redis. Это исключает дублирование данных и единственную точку рассогласования.

Структура JSON-записи заглушки описана в листинге~\ref{lst:mock-record}.

\begin{listing}[H]
\caption{Структура JSON-записи заглушки в Redis}
\label{lst:mock-record}
\begin{minted}{json}
{
    "hostname":           "test-server-01",
    "port":               8105,
    "pid":                24731,
    "status":             "RUNNING",
    "jvm_args":           "-Xms128m -Xmx256m",
    "rate_limit":         500,
    "rate_limit_enabled": false,
    "registered_at":      "2024-11-01T10:30:00Z",
    "started_at":         "2024-11-01T10:35:00Z"
}
\end{minted}
\end{listing}

Ключ: \texttt{mocks:\{filename\}}, тип: String (JSON).

Описание полей записи заглушки приведено в таблице~\ref{tab:mock-fields}.

% \begin{table}[H]
% \caption{Поля JSON-записи заглушки}
% \label{tab:mock-fields}
% \centering
% \begin{tabular}{|p{3.5cm}|p{2.5cm}|p{7cm}|}
% \hline
% Поле & Тип & Описание \\
% \hline
% \texttt{hostname} & String & Hostname целевого хоста, на котором запущена заглушка \\
% \hline
% \texttt{port} & Integer / null & TCP-порт, на котором слушает процесс заглушки \\
% \hline
% \texttt{pid} & Integer / null & Идентификатор процесса на целевом хосте \\
% \hline
% \texttt{status} & String & Текущий статус: \texttt{REGISTERED}, \texttt{RUNNING}, \texttt{STOPPED}, \texttt{ERROR} \\
% \hline
% \texttt{jvm\_args} & String & Аргументы JVM, заданные оператором при регистрации \\
% \hline
% \texttt{rate\_limit} & Integer & Максимальное количество запросов в секунду \\
% \hline
% \texttt{rate\_limit\_enabled} & Boolean & Флаг активности механизма Rate Limiter \\
% \hline
% \texttt{registered\_at} & ISO~8601 & Дата и время регистрации заглушки \\
% \hline
% \texttt{started\_at} & ISO~8601 / null & Дата и время последнего успешного запуска \\
% \hline
% \end{tabular}
% \end{table}

\begin{table}[H]
\caption{Поля JSON-записи заглушки}
\label{tab:mock-fields}
\centering
\begin{tblr}{
  colspec = {|l|l|X|},
  hlines,
  width = \linewidth,
}
Поле &                          Тип &               Описание \\
\texttt{hostname} &             String &            Hostname целевого хоста, на котором запущена заглушка \\
\texttt{port} &                 Integer / null &    TCP-порт, на котором слушает процесс заглушки \\
\texttt{pid} &                  Integer / null &    Идентификатор процесса на целевом хосте \\
\texttt{status} &               String &            Текущий статус: \texttt{REGISTERED}, \texttt{RUNNING}, \texttt{STOPPED}, \texttt{ERROR} \\
\texttt{jvm\_args} &            String &            Аргументы JVM, заданные оператором при регистрации \\
\texttt{rate\_limit} &          Integer &           Максимальное количество запросов в секунду \\
\texttt{rate\_limit\_enabled} & Boolean &           Флаг активности механизма Rate Limiter \\
\texttt{registered\_at} &       ISO~8601 &          Дата и время регистрации заглушки \\
\texttt{started\_at} &          ISO~8601 / null &   Дата и время последнего успешного запуска \\
\end{tblr}
\end{table}

Поля \texttt{port}, \texttt{pid} и \texttt{started\_at} заполняются только после успешного запуска процесса и сбрасываются в \texttt{null} при его остановке или удалении. Для получения полного пути к файлу при выполнении SSH-команд Lifecycle Manager читает поле \texttt{working\_dir} из записи \texttt{hosts:\{hostname\}} и конкатенирует его с именем файла.

\subsection{Структура данных хоста}

Запись о целевом хосте хранится по ключу вида \texttt{hosts:\{hostname\}}, где \texttt{hostname} --- сетевое имя или IP-адрес хоста. Hostname одновременно выполняет роль уникального идентификатора записи и адреса для SSH/SCP-подключения, что исключает необходимость в отдельном поле адреса. Прямой доступ к конфигурации хоста по известному hostname осуществляется за одну операцию чтения без дополнительного индекса.

Структура JSON-записи хоста описана в листинге~\ref{lst:host-record}.

\begin{listing}[H]
\caption{Структура JSON-записи хоста в Redis}
\label{lst:host-record}
\begin{minted}{json}
{
    "ssh_port":        22,
    "account_uuid":    "e5f6a7b8-c9d0-1234-efab-567890abcdef",
    "working_dir":     "/opt/mock-services",
    "java_path":       "/usr/lib/jvm/java-17-openjdk-amd64/bin/java",
    "mock_port_min":   8100,
    "mock_port_max":   8200,
    "description":     "Стенд нагрузочного тестирования – контур препром.",
    "status":          "AVAILABLE",
    "last_checked_at": "2024-11-01T10:29:55Z"
}
\end{minted}
\end{listing}

Ключ: \texttt{hosts:\{hostname\}}, тип: String (JSON).

Поле \texttt{status} отражает результат последней проверки доступности SSH-соединения, выполняемой фоновой задачей с настраиваемым интервалом. Допустимые значения: \texttt{AVAILABLE} (соединение успешно установлено), \texttt{UNAVAILABLE} (попытка соединения завершилась ошибкой) и \texttt{UNKNOWN} (проверка ещё не выполнялась). Поле \texttt{last\_checked\_at} фиксирует время последней проверки для отображения актуальности статуса в веб-интерфейсе.

Поле \texttt{account\_uuid} содержит UUID учётной записи SSH из реестра \texttt{accounts}, используемой для подключения к данному хосту. Lifecycle Manager читает это поле при каждой SSH/SCP-операции для получения соответствующей записи \texttt{accounts:\{uuid\}} с зашифрованным паролем.

\subsection{Структура данных учётной записи}

Запись об учётной записи ОС Linux хранится по ключу вида \texttt{accounts:\{uuid\}}. Учётная запись содержит данные, необходимые для аутентификации при SSH/SCP-операциях на целевых хостах. Пароль хранится в зашифрованном виде --- компонент Crypto Service шифрует его алгоритмом AES-128 (Fernet) перед записью в Redis. Ключ шифрования хранится на файловой системе сервера в файле с правами 600.

Структура JSON-записи учётной записи описана в листинге~\ref{lst:account-record}.

\begin{listing}[H]
\caption{Структура JSON-записи учётной записи в Redis}
\label{lst:account-record}
\begin{minted}{json}
{
    "username":     "mock-runner",
    "password_enc": "gAAAAABl...base64-encoded-ciphertext...",
    "description":  "Сервисный пользователь для запуска заглушек",
    "created_at":   "2024-10-15T09:00:00Z"
}
\end{minted}
\end{listing}

Ключ: \texttt{accounts:\{uuid\}}, тип: String (JSON).

Поле \texttt{password\_enc} содержит Base64-кодированный зашифрованный блок, полученный библиотекой \texttt{cryptography} (Fernet). При необходимости установить SSH-соединение Crypto Service читает это поле и выполняет дешифрование в оперативной памяти процесса; расшифрованный пароль никогда не записывается на диск и не сохраняется в Redis. Такой подход обеспечивает защиту паролей в случае несанкционированного получения дампа данных Redis.

\subsection{Структуры данных реестров}

Для хранения множеств идентификаторов всех зарегистрированных объектов используется тип Set. Каждый реестр хранится под фиксированным ключом и содержит строковые идентификаторы соответствующих объектов.

Реестр заглушек (\texttt{mocks:registry}) --- множество имён файлов всех зарегистрированных заглушек. При регистрации новой заглушки её имя добавляется командой SADD, при удалении --- удаляется командой SREM. Команда SMEMBERS возвращает полный список для формирования Dashboard.

Реестр хостов (\texttt{hosts:registry}) --- множество hostname-идентификаторов зарегистрированных целевых хостов (сетевое имя или IP-адрес). При регистрации нового хоста его hostname добавляется командой SADD, при удалении --- командой SREM. Команда SMEMBERS возвращает полный список для отображения в интерфейсе настроек.

Реестр учётных записей (\texttt{accounts:registry}) --- множество UUID зарегистрированных учётных записей SSH.

Применение типа Set гарантирует уникальность идентификаторов: повторное добавление уже существующего элемента командой SADD не создаёт дубликата и возвращает 0, что используется системой для обнаружения попытки регистрации заглушки с дублирующим именем файла.

ER-диаграмма приведена на рисунке~\ref{fig:redis-relations}.

\begin{figure}[H]
\centering
\includegraphics[width=0.85\textwidth]{redis-ER.png}
\caption{ER-диаграмма объектов в Redis}
\label{fig:redis-relations}
\end{figure}

\subsection{Структура данных счётчика Rate Limiter}

Алгоритм Rate Limiter реализован на основе паттерна «Счётчик с фиксированным временным окном» (Fixed Window Counter). Для каждой заглушки и каждого временного окна в Redis хранится отдельный ключ-счётчик. Структура ключа: \texttt{rate:\{filename\}:\{window\}}, где \texttt{\{window\}} --- метка временного окна, вычисляемая как целочисленное деление текущего Unix-timestamp на длину окна в секундах.

Например, при длине окна 1~секунда и текущем времени 1700000060.7 метка окна равна 1700000060, а ключ для заглушки \texttt{payment-service-stub.jar} примет вид \texttt{rate:payment-service-stub.jar:1700000060}. Все запросы, поступившие в течение одной секунды, инкрементируют один и тот же счётчик.

Алгоритм обработки каждого запроса через Rate Limiter описан в листинге~\ref{lst:rate-limiter}.

\begin{listing}[H]
\caption{Псевдокод алгоритма Fixed Window Counter}
\label{lst:rate-limiter}
\begin{minted}{python}
def check_rate_limit(filename: str, limit: int, window_size: int = 1) -> bool:
    # Шаг 1: вычислить метку текущего окна
    window = floor(current_unix_timestamp() / window_size)
    key    = f"rate:{filename}:{window}"

    # Шаг 2: атомарный инкремент счётчика
    count = INCR(key)

    # Шаг 3: установить TTL при создании нового окна
    if count == 1:
        EXPIRE(key, window_size   2)  # TTL = 2 окна (запас на сдвиг времени)

    # Шаг 4: проверить лимит
    if count > limit:
        return False  # запрос отклонён -> HTTP 429
    return True       # запрос разрешён
\end{minted}
\end{listing}

Установка TTL через команду EXPIRE выполняется только при создании нового счётчика (когда INCR возвращает 1), то есть при первом запросе в новом временном окне. TTL устанавливается равным удвоенному размеру окна для гарантированного самоочищения устаревших счётчиков даже при незначительных расхождениях системного времени между клиентом и сервером. Это исключает накопление «мёртвых» ключей в Redis и предотвращает утечку памяти при длительной эксплуатации системы.

Атомарность операции INCR является ключевым свойством данного подхода: Redis выполняет её как единую неделимую операцию, что гарантирует корректный подсчёт при конкурентном поступлении тысяч запросов в рамках одного временного окна без использования транзакций или блокировок.

\subsection{Глобальные настройки системы}

Глобальные конфигурационные параметры приложения, не привязанные к конкретному хосту или заглушке, хранятся под ключом \texttt{settings:global}. Запись содержит параметры, задаваемые при первоначальной настройке системы и редко изменяемые в процессе эксплуатации.

Структура JSON-записи глобальных настроек описана в листинге~\ref{lst:global-settings}.

\begin{listing}[H]
\caption{Структура глобальных настроек в Redis}
\label{lst:global-settings}
\begin{minted}{json}
{
    "rate_limit_window_size":  1,
    "host_check_interval_sec": 30,
    "proxy_timeout_sec":       10,
    "log_retention_lines":     1000
}
\end{minted}
\end{listing}

Ключ: \texttt{settings:global}, тип: String (JSON).

Описание полей глобальных настроек приведено в таблице~\ref{tab:global-settings}.

\begin{table}[H]
\caption{Поля глобальных настроек системы}
\label{tab:global-settings}
\centering
\begin{tabular}{|p{4.5cm}|p{2cm}|p{6.5cm}|}
\hline
Поле & Тип & Описание \\
\hline
\texttt{rate\_limit\_window\_size} & Integer & Размер временного окна алгоритма Rate Limiter в секундах \\
\hline
\texttt{host\_check\_interval\_sec} & Integer & Интервал фоновой проверки доступности хостов в секундах \\
\hline
\texttt{proxy\_timeout\_sec} & Integer & Таймаут ожидания ответа от целевой заглушки при проксировании в секундах \\
\hline
\texttt{log\_retention\_lines} & Integer & Максимальное количество строк журнала процесса, хранимых в памяти \\
\hline
\end{tabular}
\end{table}

\subsection{Персистентность данных и стратегия восстановления}

Redis настроен в режиме AOF (Append Only File) с параметром \texttt{appendfsync everysec}: каждая операция записи фиксируется в журнале, синхронизируемом с диском раз в секунду. Максимальная потеря данных при аварии составляет не более одной секунды.

Стратегия восстановления после перезапуска: управляющее приложение читает реестры из Redis и для каждой заглушки со статусом \texttt{RUNNING} проверяет фактическое наличие процесса на целевом хосте через SSH. Если процесс обнаружен --- статус подтверждается. Если нет --- статус обновляется до \texttt{STOPPED}, поля \texttt{port}, \texttt{pid} и \texttt{started\_at} сбрасываются. Бинарный файл на целевом хосте не затрагивается и остаётся в рабочей директории, готовый к повторному запуску.

В таблице~\ref{tab:redis-ops} представлена сводная матрица операций чтения и записи Redis по компонентам системы.

\begin{table}[H]
\caption{Матрица операций Redis по компонентам системы}
\label{tab:redis-ops}
\centering
\begin{tabular}{|p{3.5cm}|p{3cm}|p{6.5cm}|}
\hline
Компонент & Операции Redis & Назначение \\
\hline
Lifecycle Manager & GET, SET, SADD, SREM, DEL & Чтение конфигураций хостов/заглушек, обновление статусов и PID \\
\hline
Proxy Engine & GET & Чтение адреса, порта и флага \texttt{rate\_limit\_enabled} заглушки \\
\hline
Rate Limiter & INCR, EXPIRE, GET & Инкремент и проверка счётчика запросов \\
\hline
Settings Module & GET, SET, SADD, SREM, WATCH, MULTI, EXEC & Управление конфигурациями хостов и учётных записей \\
\hline
Metrics Exporter & SMEMBERS, GET (Pipeline) & Агрегация метрик по всем заглушкам и хостам \\
\hline
\end{tabular}
\end{table}

Для операции Metrics Exporter применяется механизм конвейеризации команд Redis (Pipeline): сначала читается полный список ключей из реестра, затем все операции чтения конфигураций отправляются единым пакетом. Это снижает суммарные сетевые задержки при агрегации данных о большом числе заглушек с~$O(n)$ до фактически~$O(1)$ по числу сетевых round-trip.

\section {Выводы по второй главе}

Во второй главе выпускной квалификационной работы проведено проектирование архитектуры и структур данных разрабатываемого программного комплекса.

На основании сравнительного анализа архитектурных подходов обоснован выбор монолитной архитектуры с модульной внутренней организацией. Показано, что данное решение оптимально соответствует требованиям проекта: минимизирует накладные расходы на проксирование, исключает зависимость от внешних оркестраторов и обеспечивает простоту эксплуатации в закрытом корпоративном контуре.

Разработана топология развёртывания системы, обеспечивающая гибкое размещение имитационных сервисов: единый управляющий процесс координирует как локальные, так и удалённые хосты по единому алгоритму через SSH/SCP без установки агентов на целевые серверы.

Архитектура системы формализована в нотации C4~Model на трёх уровнях детализации: контекстном, контейнерном и компонентном. Определены шесть программных компонентов --- Lifecycle Manager, Proxy Engine, Rate Limiter, Settings Module, Crypto Service и Metrics Exporter --- и описаны типовые сценарии их взаимодействия.

Спроектированы структуры данных Redis, обеспечивающие хранение конфигураций заглушек, хостов и учётных записей в формате JSON с доступом за одну операцию GET. Реализован алгоритм Fixed Window Counter для ограничения интенсивности запросов с автоматической очисткой устаревших счётчиков через механизм TTL. Настроена персистентность данных в режиме AOF с разработанной стратегией восстановления состояний после перезапуска системы.